% 为本书第4册（大学化学，第三学年）的照片与人工智能插图列出出处。机器可读的许可记录位于 images/book4/CREDITS.md；两个文件须保持同步。本书的编写者负责维护此文件：照片入库后，在第一段之后的 itemize 中为每张照片添加一个 \item。
\chapter*{Image Credits}

本书中的图示均为原创。照片（如有使用）依其各自的自由许可使用，谨向其作者致谢；每张照片完整的来源链接与许可网址列于本书代码库中的 \texttt{images/book4/CREDITS.md}。
危险象形图采用《全球化学品统一分类和标签制度》（Globally Harmonized System of Classification and Labelling of Chemicals）的图样，取自 \TeX\ Live 中 \texttt{ghsystem} 宏包所提供的版本。

\bigskip
除照片和原创图示外，部分图为人工智能生成的插图，由 OpenAI 图像生成功能依据本书编辑撰写的提示词制作，并在收录前经过化学准确性审校。每幅生成图像的完整提示词记录于代码库中的 \texttt{images/book4/ai/PROMPTS.md}。
\bigskip
\noindent 照片，按章列出：
\begin{itemize}\small
  \item 第~1~章：埃尔温·薛定谔，1933 年，诺贝尔基金会，公有领域（维基共享资源）。
  \item 第~4~章：雪晶，Wilson A. Bentley（约 1902 年），公有领域（维基共享资源）。
  \item 第~6~章：高原上的无线电天线，ESO/B.~Tafreshi（twanight.org），CC BY 4.0（维基共享资源）。
  \item 第~7~章：日光下与紫外光下的萤石，由 Masha Milshina 拍摄，CC BY 4.0（维基共享资源）。
  \item 第~8~章：一间核磁共振实验室，由 Shandchem 拍摄，CC BY 2.0（维基共享资源）。
  \item 第~9~章：威廉·劳伦斯·布拉格，诺贝尔基金会，公有领域（维基共享资源）。
  \item 第~10~章：路德维希·玻尔兹曼的墓碑，由 Feldkurat Katz 拍摄，CC BY-SA 4.0（维基共享资源）。
  \item 第~13~章：Belousov--Zhabotinsky 螺旋图样，由 Stephen Morris 拍摄，CC BY 2.0（维基共享资源）。
  \item 第~14~章：南极臭氧空洞，NASA，公有领域（维基共享资源）。
  \item 第~16~章：催化转化器蜂窝结构的电子显微照片，由 Galadrid 拍摄，CC BY-SA 4.0（维基共享资源）。
  \item 第~17~章：阳光束与丁达尔效应，由 Kevin c higgins 拍摄，CC BY-SA 4.0（维基共享资源）。
  \item 第~18~章：红宝石与祖母绿晶体，由 Robert M. Lavinsky 拍摄，CC BY-SA 3.0（维基共享资源）。
  \item 第~20~章：二茂铁晶体，由 Andrea Sella 拍摄，CC BY 4.0（维基共享资源）。
  \item 第~22~章：单晶硅锭，由 Sebastian Wallroth 拍摄，公有领域（维基共享资源）。
  \item 第~23~章：硅气凝胶上的一朵花，NASA/JPL，公有领域；莱克格斯杯，由 Marie-Lan Nguyen 拍摄，CC BY 2.5（维基共享资源）。
  \item 第~24~章：大西洋鲎，由 Kaldari 拍摄，CC0（维基共享资源）。
  \item 第~27~章：除虫菊雏菊，由 Soramimi 拍摄，CC BY-SA 4.0（维基共享资源）。
  \item 第~30~章：太平洋紫杉的树皮，由 JOE BLOWE（theforestprimeval）拍摄，CC BY-SA 2.0（维基共享资源）。
  \item 第~32~章：被 Landsat 8 卫星拍到的藻华，USGS 与 NASA，公有领域。
  \item 第~33~章：双真空与惰性气体总管装置，由 Polimerek 拍摄，CC BY-SA 3.0（维基共享资源）。
\end{itemize}
\clearpage
